\par\medskip
\hypertarget{native-073-article055}{}\pdfbookmark[2]{55.}{article-055}
\begin{center}55.\end{center}
Wir bezeichnen wie oben die Länge des aufsteigenden Knotens der Bahn auf der Ekliptik mit $\AscendingNode$ und die Neigung mit $i$. Ferner bezeichnen wir die Länge des aufsteigenden Knotens der neuen Ebene auf der Ekliptik mit $n$ und ihre Neigung mit $\epsilon$. Den Abstand des aufsteigenden Knotens der Bahn auf der neuen Ebene vom aufsteigenden Knoten der neuen Ebene auf der Ekliptik (den Bogen $n\AscendingNode'$ in Abb. 2) bezeichnen wir mit $\AscendingNode'$, die Neigung der Bahn gegen die neue Ebene mit $i'$. Den Bogen von $\AscendingNode$ nach $\AscendingNode'$ in Bewegungsrichtung bezeichnen wir schließlich mit $\Delta$. Unser sphärisches Dreieck hat dann die Seiten $\AscendingNode-n$, $\AscendingNode'$ und $\Delta$ und die gegenüberliegenden Winkel $i'$, $180^0-i$ und $\epsilon$. Nach den Formeln des vorigen Artikels ergibt sich
\par\Needspace{11\baselineskip}
\SourceLeaf{074}
\[
\sin\tfrac{1}{2}i'\sin\tfrac{1}{2}(\AscendingNode'+\Delta) = \sin\tfrac{1}{2}(\AscendingNode-n)\sin\tfrac{1}{2}(i+\epsilon)
\]
\[
\sin\tfrac{1}{2}i'\cos\tfrac{1}{2}(\AscendingNode'+\Delta) = \cos\tfrac{1}{2}(\AscendingNode-n)\sin\tfrac{1}{2}(i-\epsilon)
\]
\[
\cos\tfrac{1}{2}i'\sin\tfrac{1}{2}(\AscendingNode'-\Delta) = \sin\tfrac{1}{2}(\AscendingNode-n)\cos\tfrac{1}{2}(i+\epsilon)
\]
\[
\cos\tfrac{1}{2}i'\cos\tfrac{1}{2}(\AscendingNode'-\Delta) = \cos\tfrac{1}{2}(\AscendingNode-n)\cos\tfrac{1}{2}(i-\epsilon).
\]
Die ersten beiden Gleichungen liefern $\frac{1}{2}(\AscendingNode'+\Delta)$ und
$\sin\frac{1}{2}i'$. Die beiden übrigen liefern $\frac{1}{2}(\AscendingNode'-\Delta)$ und $\cos\frac{1}{2}i'$.
Aus $\frac{1}{2}(\AscendingNode'+\Delta)$ und $\frac{1}{2}(\AscendingNode'-\Delta)$ erhält man $\AscendingNode'$ und
$\Delta$. Aus $\sin\frac{1}{2}i'$ oder $\cos\frac{1}{2}i'$ bestimmt man
$i'$ (die Übereinstimmung beider Ergebnisse dient zur Kontrolle der Rechnung). Die Ungewissheit, ob $\frac{1}{2}(\AscendingNode'+\Delta)$ und
$\frac{1}{2}(\AscendingNode'-\Delta)$ zwischen $0$ und $180^0$ oder zwischen $180^0$ und $360^0$ zu wählen sind,
beseitigt man, indem man sie so wählt, dass sowohl $\sin\frac{1}{2}i'$ als auch $\cos\frac{1}{2}i'$ positiv werden,
da $i'$ hier seiner Bedeutung nach unter $180^0$ liegen muss.
